\input{../preamble.tex}
\SetDate[25/03/2026]
\setbeamertemplate{page number in head/foot}{}

\begin{document}

\section{Automatismes}

\begin{frame}

\centering \huge
Automatismes

\large
Calculatrice interdite

\end{frame}

\subsection{Univers probabilités}



\framedelayed{
	On lance un D5 bien équilibré, dé à 5 faces numérotées de 1 à 5.
	On regarde le numéro de la face du dessus.
	\boxAB{
		Donner la probabilité d'obtenir 4.
	}{
		Donner l'univers $\Omega$ de l'expérience.
	}
}{
	On lance un D5 bien équilibré, dé à 5 faces numérotées de 1 à 5.
	On regarde le numéro de la face du dessus.
	\boxAB{
		$P(4) = \dfrac15$.
	}{
		$\Omega = \bigset{1 ; 2 ; 3 ; 4 ; 5}$.
	}
}


\framedelayed{
	\vspace{20pt}
	On tire au hasard une boule dans une urne qui contient
	\begin{itemize}	
		\item 2 boules rouges ;
		\item 1 boule noire ;
		\item 3 boules vertes ;
	\end{itemize}
	et on regarde sa couleur.
	\boxAB{
		Donner l'univers $\Omega$ de l'expérience.
	}{
		Donner la probabilité d'obtenir une boule rouge.
	}
}{
	\vspace{20pt}
	On tire au hasard une boule dans une urne qui contient
	\begin{itemize}	
		\item 2 boules rouges ;
		\item 1 boule noire ;
		\item 3 boules vertes ;
	\end{itemize}
	et on regarde sa couleur.
	\boxAB{
		$\Omega = \bigset{ \text{Rouge} ; \text{Noir} ; \text{Vert} }$.
	}{
		$P(\text{Rouge}) = \dfrac26 = \dfrac13$.
	}
}

\subsection{Union et intersection}

\framedelayed{
	Considérons les ensembles $A$ et $B$ suivants.
		\begin{align*}
			A &= \bigset{ 4 ; 6 ; 41 ; 36 ; 7 } \\
			B &= \bigset{ 3 ; 7 ; 6 ; 1 ; 36}
		\end{align*}
	\boxAB{
		Donner $A \cap B$.
	}{
		Donner $A \cup B$.
	}
}{
	\begin{align*}
		A &= \bigset{ 4 ; 6 ; 41 ; 36 ; 7 } \\
		B &= \bigset{ 3 ; 7 ; 6 ; 1 ; 36}
	\end{align*}
	\boxAB{
		$A \cap B = \bigset{ 6 ; 36 ; 7}$
	}{
		$A \cup B = \bigset{ 4 ; 6 ; 41 ; 36 ; 7 ; 3 ; 1}$
	}
}

\framedelayed{
	Considérons les ensembles $A$ et $B$ suivants.
		\begin{align*}
			A &= \bigset{ 5 ; 8 ; 1 ; 85 ; 2 } \\
			B &= \bigset{ 5 ; 2 ; 1 ; 9 ; 11}
		\end{align*}
	\boxAB{
		Donner $A \cup B$.
	}{
		Donner $A \cap B$.
	}
}{
		\begin{align*}
			A &= \bigset{ 5 ; 8 ; 1 ; 85 ; 2 } \\
			B &= \bigset{ 5 ; 2 ; 1 ; 9 ; 11}
		\end{align*}
	\boxAB{
		$A \cup B = \bigset{ 5 ; 8 ; 1 ; 85 ; 2 ; 9 ; 11 }$
	}{
		$A \cap B = \bigset{ 5 ; 1 ; 2}$
	}
}

%\subsection{Courbe représentative}
%
%\framedelayed{
%	Soient 
%		\begin{align*}
%			f(x) = 2x+1 && \et && g(x) = x+2
%		\end{align*}
%	deux fonctions sur $\R$.
%	\boxAB{
%		Donner un point appartenant à $\C_f$.
%	}{
%		Donner un point appartenant à $\C_g$.
%	}
%}{
%	Soient 
%		\begin{align*}
%			f(x) = 2x+1 && \et && g(x) = x+2
%		\end{align*}
%	deux fonctions sur $\R$.
%	\boxAB{
%		$f(0) = 1$ donc $(0 ; 1)$ appartient à $\C_f$.
%	}{
%		$g(12) = 14$ donc $(12 ; 14)$ appartient à $\C_g$.
%	}
%}


\section{Solutions}

\shipoutAnswer

\end{document}